\section{Frames along geodesics and Hessian comparison}\label{sec:frames}

Let $(M,d,m)$ be a full-support noncollapsed $\RCD(K,n)$ space,
where $m=\mathcal H^n$ and $n\geq2$ is an integer.
Subsections~\ref{frames:sub-countable}--\ref{frames:sub-paths} use only
these ambient assumptions.  Their path calculus and trial approximation
apply to optimal geodesic plans with bounded positions, bounded speed,
and uniformly bounded compression throughout the interval in question.
The sectional tensor lower bound enters the radial application in
Subsection~\ref{frames:sub-radial} through Theorem~\ref{heat:index}.

Theorem~\ref{heat:index} bounds the distributional Hessian by the
energies of global trials.  To insert Jacobi trials, we need frames at
every time on each positive-time interval along almost every calibrated
geodesic.  We construct a countable family of test gradients whose rank
defect is avoided by almost every entire path, recover parallel transport by
finite-dimensional ODEs, and obtain one countable catalog of global
trials that is dense separately along almost every path.  The common
terminal-point prescription then lets us take the infimum of their
energy functions at each terminal point, yielding the measure bound in
Theorem~\ref{frames:hessian}.

Write $\mathcal R_n$ for the set of points all of whose measured tangents
are the normalized Euclidean $n$-space.  Noncollapse implies
$m(M\setminus\mathcal R_n)=0$ and that the tangent module has dimension
$n$ almost everywhere.  We use the local volume estimate
\eqref{eq:local-volume} and the local Poincar\'e inequality, with
constants uniform on bounded regions.

A compact test function means an element of $\Test(M)$ with compact
support.  A compact test gradient is the gradient of such a function.
For locally Lipschitz $h$, we write $H_h=\Hess h$ for the
distributional functional \eqref{eq:weak-hessian}.  For a local
potential, weights are compactly supported in its domain; locality
makes the evaluation independent of an extension of $h$.

\subsection{A countable family of frames outside small exceptional sets}
\label{frames:sub-countable}

A frame defined almost everywhere may still fail at a time visited by
a path.  We therefore control the size of the set of regular points
where every selected frame fails.  The following proposition separates
a polar set, needed to choose common Sobolev representatives, from a
rank defect of vanishing $(n-1)$-dimensional Hausdorff measure.
Lemma~\ref{frames:path-avoidance} will show that almost every entire
bounded-speed, bounded-compression path avoids both sets.

\begin{proposition}[Countable harmonic frames]
 \label{frames:countable-frames}
 There is a countable family $\{g_j\}_{j\geq1}$ of compact test
 functions, a set $N$ of $2$-capacity zero, and a Borel set
 $S\subset\mathcal R_n\setminus N$ with
 $\mathcal H^{n-1}(S)=0$ such that the following holds.
 At every $x\in\mathcal R_n\setminus(N\cup S)$, some $n$ members of
 $\{\nabla g_j\}$ have positive definite Gram matrix, evaluated in
 common precise Sobolev representatives.
 All finite Gram matrices are positive semidefinite and have rank
 at most $n$ outside $N$.
 Any prescribed countable family of compact test functions may be
 included among the $g_j$.
\end{proposition}

The proof uses two content estimates.  The first turns small local
content into a Hausdorff-null set; the second bounds the content of a
Gram-matrix rank defect by its Sobolev energy.  Hausdorff content is
defined using sums of powers of diameters; changes between radii and
diameters will be absorbed in constants.

\begin{lemma}[Small local content]\label{frames:content-null}
 Let $E$ be a subset of a separable metric space and let $s>0$.  If
 \begin{equation}\label{frames:vanishing-content}
  \mathcal H^s_\infty(E\cap B(x,r))=o(r^s)
  \quad\text{as }r\downarrow0\quad\text{for every }x\in E,
 \end{equation}
 then $\mathcal H^s(E)=0$.  It suffices to assume
 \eqref{frames:vanishing-content} along all sufficiently small dyadic
 radii.
\end{lemma}

\begin{proof}
 Bounds along dyadic radii imply bounds at every radius by enlarging
 to a dyadic radius at most twice as large.  Fix $\varepsilon>0$, to be
 chosen sufficiently small in terms of $s$, and put
 \[
 E_j=\{x\in E:\mathcal H^s_\infty(E\cap B(x,r))
                \leq\varepsilon r^s\text{ for every }0<r<1/j\}.
 \]
 The sets $E_j$ cover $E$.  Fix $x\in E_j$ and $R<1/j$.
 A content cover of $E_j\cap B(x,R)$ can be recentered in $E_j$:
 discard covering sets not meeting this set and replace a covering set
 of positive diameter $d_i$ by a ball of radius $2d_i$ centered at a
 point of its intersection with $E_j$.  Replace singleton covering
 sets by balls with arbitrarily small positive total $s$-cost.
 Choosing the original cover within an arbitrarily small error of
 the content gives child balls $B(x_i,r_i)$ with
 \[
  x_i\in E_j,\qquad \sum_i r_i^s\leq\theta R^s,
  \qquad 0<\theta<1,
 \]
 provided $\varepsilon$ was chosen small enough.  In particular
 $r_i\leq\theta^{1/s}R<1/j$.

 Repeat this procedure on $E_j\cap B(x_i,r_i)$, which may enlarge
 the part of the parent set being covered but still covers it.
 At generation $k$ the total diameter cost is at most
 $2^s\theta^kR^s$, and every diameter is at most
 $2\theta^{k/s}R$.  Both bounds tend to zero.  Thus
 $\mathcal H^s(E_j\cap B(x,R))=0$.
 Separability supplies countably many such root balls covering $E_j$.
 Taking the countable union over $j$ proves the claim.  No finiteness
 assumption on $\mathcal H^s(E)$ is used.
\end{proof}

\begin{lemma}[Energy controls the content of a rank defect]
 \label{frames:energy-content}
 Work in a bounded region where \eqref{eq:local-volume} and the
 Poincar\'e inequality hold with fixed constants.  There are an
 enlargement factor $L$ and $\varepsilon_0>0$ with the following
 property.  Let $P$ be a bounded Sobolev matrix on $B(x,Lr)$ and suppose
 \[
  \int_{B(x,Lr)}|P-I|\dd m\leq\varepsilon_0r^n.
 \]
 For $n-2<s$ let $E\subset B(x,r)$ consist of precise-representative
 points of $P$ at which $|P-I|\geq1$.  Then
 \begin{equation}\label{frames:energy-content-bound}
  \mathcal H^s_\infty(E)
   \leq C_s r^{s+2-n}
             \int_{B(x,Lr)}|\nabla P|^2\dd m.
 \end{equation}
 The constants depend only on $s,n$ and the local volume and
 Poincar\'e constants.
\end{lemma}

\begin{proof}
 The volume lower bound makes the average of $P$ on $B(y,r)$ within
 $1/4$ of $I$, simultaneously for all $y\in B(x,r)$, if
 $\varepsilon_0$ is small enough.  Let $\lambda$ be the enlargement
 factor in the Poincar\'e inequality, and choose $L>2+5\lambda$.
 For $t_j=2^{-j}r$, the difference between the averages on
 $B(y,t_j)$ and $B(y,t_{j+1})$ is at most
 \[
  C t_j^{1-n/2}
       \left(\int_{B(y,\lambda t_j)}|\nabla P|^2\dd m\right)^{1/2}.
 \]
 If, for every $j\geq0$,
 \[
  \int_{B(y,\lambda t_j)}|\nabla P|^2\dd m
       <c_0t_j^s r^{n-2-s},
 \]
 the sum of these differences is bounded by
 \[
  C\sqrt{c_0}\sum_{j\geq0}2^{-j(s-n+2)/2}.
 \]
 Because $s>n-2$, this is less than $1/2$ for an appropriate $c_0>0$.
 The averages converge to the precise value at $y$, so this
 contradicts $|P(y)-I|\geq1$.
 Consequently, each $y\in E$ has a radius $t_y\leq r$ for which
 \begin{equation}\label{frames:energy-threshold}
  \int_{B(y,\lambda t_y)}|\nabla P|^2\dd m
       \geq c_0t_y^s r^{n-2-s}.
 \end{equation}
 The $5r$ covering lemma selects countably many disjoint balls
 $B(y_i,\lambda t_i)$ whose fivefold enlargements cover $E$.
 They lie in $B(x,Lr)$.  Summing
 \eqref{frames:energy-threshold} over these disjoint balls bounds
 $\sum_i t_i^s$ by the right-hand side of
 \eqref{frames:energy-content-bound}, up to a fixed factor.
 Their fivefold enlargements give the required content cover.
\end{proof}

\begin{proof}[Proof of Proposition~\ref{frames:countable-frames}]
 \textbf{Selection at all dyadic scales.}
 Put $r_k=2^{-k}$, and choose a countable $r_k$-net $N_k$ in $M$.
 For each integer $A\geq4$, each $y\in N_k$, and each integer $q\geq1$,
 select one harmonic $n$-tuple on $B(y,2Ar_k)$ with splitting error
 at most $1/(n^2q)$, whenever such a tuple exists.  The splitting-map
 estimates give a dimensional gradient bound and
 \begin{align*}
  \frac{1}{m(B(y,2Ar_k))}\int_{B(y,2Ar_k)}|P-I|\dd m&\leq q^{-1},\\
  \frac{(2Ar_k)^2}{m(B(y,2Ar_k))}
          \int_{B(y,2Ar_k)}|\Hess u^a|^2\dd m&\leq q^{-1},
 \end{align*}
 for its components $u^a$ and their Gram matrix $P$.
 The existence input is the harmonic splitting-map theorem near a
 Euclidean rescaling; see \cite[Proposition~1.4]{BPS2021}.
 To match that theorem's product formulation, use the auxiliary
 synthetic dimension \(N=n+1\), Euclidean factor dimension \(n\),
 and a point as the \(\RCD(0,1)\) residual factor. Its radius is
 \(2Ar_k/5\), so its output ball is \(B(y,2Ar_k)\). The rescaled
 curvature bound tends to zero. This auxiliary parameter does not
 change the actual dimension \(n\), the measure \(\mathcal H^n\),
 or the volume estimates used here.

 Subtract the values $u^a(y)$ and multiply by a good test cutoff
 supported in $B(y,2Ar_k)$ and equal to one on $B(y,Ar_k)$.
 These are compact test functions.  Indeed, on the support of the
 cutoff $\eta$, the functions and their gradients are bounded and
 their Hessians are locally in $L^2$, and
 \[
  \Delta(\eta u^a)=u^a\Delta\eta
                       +2\langle\nabla\eta,\nabla u^a\rangle.
 \]
 The right-hand side is in $W^{1,2}$ by the product rule and the
 local Hessian bounds.  The construction therefore changes none of
 the Gram entries on the inner ball.  The collection of all selected
 components is countable.  Adjoin the prescribed countable family.

 Fix a regular point $x$, a bounded region containing it, and an
 integer $A>L+1$, with $L$ as in
 Lemma~\ref{frames:energy-content}.  Choose $y_k\in N_k$ with
 $d(x,y_k)\leq r_k$.  The rescalings at $y_k$ converge to Euclidean
 space at each fixed enlargement of $r_k$.  To check this, take any
 subsequence of rescalings at $x$.  Their limit is Euclidean by
 regularity.  The moving basepoints are at bounded rescaled
 distance and hence have a further convergent subsequence; the
 resulting Euclidean pointed space is isometric to the one pointed
 at the origin.  This proves the assertion for the original sequence
 as well.  The same reasoning applies to normalized measures.

 It follows that, for each fixed $q$, the selected harmonic tuple
 with threshold $q$ is available at $(y_k,2Ar_k)$ for every
 sufficiently large $k$.  Choose $q_k\leq k$ maximal among the
 available thresholds, ignoring the finitely many initial scales
 where there is no choice.  Then $q_k\to\infty$.
 On $B(x,Lr_k)$, the resulting Gram matrix $P_k$ agrees with that
 of the harmonic tuple.  Metric compatibility, bounded gradients,
 and \eqref{eq:local-volume} imply
 \begin{equation}\label{frames:all-dyadic-smallness}
  \int_{B(x,Lr_k)}
       \bigl(|P_k-I|+r_k^2|\nabla P_k|^2\bigr)\dd m
        \leq C_Aq_k^{-1}r_k^n=o(r_k^n).
 \end{equation}
 Thus the conclusion holds at every sufficiently small dyadic scale,
 rather than only along a sparse subsequence.

 \textbf{Common precise representatives.}
 Every pairwise inner product of the countably many gradients is a
 bounded Sobolev function.  On the present complete, locally doubling
 Poincar\'e space, such functions have $L^2$-Lebesgue representatives,
 given by limits of ball averages outside a set of $2$-capacity zero,
 and these are their quasi-continuous representatives; see
 \cite[Propositions~4.8--4.9]{BB2018}. The equal-norm identification
 of Newtonian and RCD weak-gradient Sobolev spaces follows from
 \cite[Theorem~7.4 and Section~8.2]{AGSWeakGradients2013}; see also
 \cite{Gigli2018}.
 Completeness makes $M$ its own completion, so the path-openness
 hypothesis of Proposition~4.9 and the measurable-inclusion
 qualification in the appended correction are automatic here.
 To identify the polar sets, let $C_W$ be the outer Sobolev capacity
 and $C_N$ the Newtonian capacity, using the same norm normalization.
 On an open set $U$, a Newtonian competitor is a Sobolev competitor
 of the same norm.  Conversely, truncate a Sobolev competitor to
 $[0,1]$ and take its equal-norm Newtonian representative.  This
 representative equals one at every one of its $L^2$-Lebesgue
 points in $U$.  The exception is $C_N$-polar, so resetting it to
 one there preserves its Newtonian class and norm.  Thus the two
 capacities agree on open sets.  Their outer regularity identifies
 their polar sets, and uniqueness identifies the resulting
 quasi-continuous representatives.
 Remove one common polar set for this countable collection.
 The rational quadratic forms of the Gram matrices are nonnegative
 almost everywhere, and all their $(n+1)$-minors vanish almost
 everywhere.  On the bounded range of any finite list of Gram
 entries, each such polynomial is Lipschitz.  Their $L^2$-mean
 convergence at a common Lebesgue point therefore implies that
 the polynomial averages converge to the polynomial of the
 precise values.  This passes the almost-everywhere positivity
 and minor identities to these values.  The polynomials are
 themselves bounded Sobolev functions.  Removing a
 further countable union of polar sets makes these identities hold
 simultaneously in the precise representatives.  This proves the
 positivity and rank assertions outside a single polar set $N$,
 which may be taken Borel.

 \textbf{Size of the rank defect.}
 Let $S$ be the set of regular points outside $N$ where no $n$-tuple
 is positive definite.  The representatives and $\mathcal R_n$ may
 be taken Borel, so $S$ is Borel.  Each chosen $n$-tuple is singular
 at every point of $S$, and hence $|P_k-I|\geq1$ there.
 Apply Lemma~\ref{frames:energy-content} and
 \eqref{frames:all-dyadic-smallness}.  For every regular $x$ and
 every $s>n-2$ this gives
 \[
  \mathcal H^s_\infty(S\cap B(x,r_k))=o(r_k^s).
 \]
 The constants are uniform on each bounded region.
 Lemma~\ref{frames:content-null}, followed by a bounded exhaustion, gives
 $\mathcal H^s(S)=0$.  In particular we take $s=n-1$.
 When $n=2$ this is $s=1$; the telescoping series in
 Lemma~\ref{frames:energy-content} has ratio $2^{-1/2}$, so no
 borderline exponent is used.
\end{proof}

\subsection{Representatives and frames on entire paths}
\label{frames:sub-paths}

For a finite measure $\pi$ on paths, write $e_t(\gamma)=\gamma_t$.
All statements about test plans below also apply to finite plans by
normalization.  Bounded positions mean that the union of the images
of paths in the support lies in a bounded subset of $M$.

\begin{lemma}[Entire-path avoidance and representatives]
 \label{frames:path-avoidance}
 Let $a<b$ and let $\pi$ be a finite plan concentrated on
 $L$-Lipschitz paths with bounded positions.  Suppose
 $(e_t)_\#\pi\leq C m$ for every $t\in[a,b]$.
 For each fixed polar set $N$ and each fixed $\mathcal H^{n-1}$-null
 set $S$, almost every path satisfies
 $\gamma([a,b])\cap(N\cup S)=\varnothing$.
 For any countable family of
 Sobolev functions, their quasi-continuous representatives agree,
 along almost every entire path, with the absolutely continuous
 representatives furnished by Sobolev differentiation along plans.
\end{lemma}

\begin{proof}
 Put $T=b-a$ and $M_\pi=\pi(C([a,b],M))$.  Sobolev differentiation
 along bounded-compression plans and the elementary one-dimensional
 bound on the supremum of an absolutely continuous function give
 \begin{equation}\label{frames:path-sup}
  \int\sup_{t\in[a,b]}|u(\gamma_t)|\dd\pi
   \leq\sqrt{CM_\pi}
            \bigl(\|u\|_{L^2(m)}+LT\|\nabla u\|_{L^2(m)}\bigr),
 \end{equation}
 where on the left we initially use the continuous path representative.
 The chain rule and its representative formulation are part of the
 RCD Sobolev calculus; see \cite{Gigli2018,Savare2014}.

 If $N$ is polar, choose open neighborhoods $U_j\supset N$ and
 $u_j\in W^{1,2}(M)$ with $u_j\geq1$ almost everywhere on $U_j$
 and $\|u_j\|_{W^{1,2}}\to0$.
 For almost every path, membership in the exceptional null subset
 of $U_j$ occupies zero time, by bounded compression and Fubini.
 A path hitting $U_j$ has a nontrivial one-sided time interval in
 $U_j$.  Its continuous representative of $u_j$ therefore has
 supremum at least one.  Estimate~\eqref{frames:path-sup} shows that
 the measure of paths hitting $N$ is zero.

 For the Hausdorff assertion set $L_*=\max\{1,L\}$.
 A path hitting $B(z,r)$ spends at least $r/L_*$ units of time in
 $B(z,2r)$ whenever $r<L_*T$.  This remains true for a hit at an
 endpoint, using a one-sided interval.  Consequently
 \[
  \pi\{\gamma:\gamma([a,b])\cap B(z,r)\ne\varnothing\}
       \leq\frac{CL_*T}{r}m(B(z,2r))\leq C' r^{n-1}.
 \]
 The last constant is uniform on the bounded region under
 consideration.  Covering an $\mathcal H^{n-1}$-null set by balls
 with arbitrarily small total $(n-1)$-cost proves its avoidance.

 Finally, approximate each of the specified Sobolev functions
 sufficiently rapidly in $W^{1,2}$ by Lipschitz functions.
 A subsequence converges to its quasi-continuous representative
 outside a polar set.  Estimate~\eqref{frames:path-sup}, applied
 to successive differences, makes the approximations converge
 uniformly along almost every path.  Those paths avoid the polar
 set just obtained; their uniform limit is therefore the
 quasi-continuous representative at every time.  It also equals
 the Sobolev path representative, since they agree for almost
 every time.  Take the countable intersection of the resulting
 conull sets.
\end{proof}

\begin{lemma}[Regularity at every time]\label{frames:regular-paths}
 Let $\pi$ be a bounded-position geodesic plan on $[a,b]$ whose
 marginals have uniformly bounded densities.  Then almost every
 path lies in $\mathcal R_n$ for every $t\in[a,b]$.
\end{lemma}

\begin{proof}
 Almost every path is regular at the endpoints and at every time
 in a fixed countable dense subset of $(a,b)$.
 Consider one such nonconstant path and fix an interior time $t$.
 Take any sequence of rescalings converging to a tangent $Y$ at
 $\gamma_t$.  At a nearby regular time $t_j$, the same rescalings
 converge to Euclidean $n$-space.  Deng's estimate, valid along
 every geodesic, implies for each fixed radius $R$ that
 \[
  d_{\mathrm{pGH}}\bigl((B_R(Y),o_Y),
                   (B_R(\mathbb R^n),0)\bigr)
       \leq C R|t-t_j|^\alpha.
 \]
 The constants can be kept fixed while $t_j$ stays in a compact
 interior interval containing $t$; see
 \cite[Theorems~1.1--1.2]{Deng2020}.
 Letting regular $t_j$ tend to $t$ proves that every pointed ball
 in $Y$ is Euclidean.  The tangent sequence was arbitrary, so
 $\gamma_t\in\mathcal R_n$.  In the noncollapsed setting the
 normalized tangent measures are the corresponding Euclidean
 measures.  Constant paths are already covered by regularity at
 one fixed time.  Endpoint regularity came directly from bounded
 compression, not from the interior estimate.
\end{proof}

Together, the preceding lemmas give a full-rank Gram matrix at every
time along almost every relevant path.  We next use its continuous
entries to define fibers and a connection, and then identify the
geodesic velocity and its endpoint traces.

\begin{proposition}[Finite-dimensional calculus along a path]
 \label{frames:path-calculus}
 Fix a countable family as in Proposition~\ref{frames:countable-frames}.
 Let $\pi$ be an optimal geodesic plan on $[a,b]$ with bounded
 positions, bounded speed, and uniformly bounded-compression
 marginals.  Along almost every path there are $n$-dimensional
 fibers at every time and local frames formed from the specified
 test gradients, with the following properties.
 \begin{enumerate}
  \item Gram entries belong to $W^{1,2}([a,b])$, are continuous,
  and differentiate by metric compatibility with $\nabla_v$.
  \item In a local frame, fields with $W^{1,2}$ scalar
  coefficients have a well-defined metric-compatible derivative
  \begin{equation}\label{frames:curve-derivative}
   \mathcal D\left(\sum_i a_i e_i\right)
        =\sum_i\bigl(a_i'e_i+a_i\nabla_v e_i\bigr).
  \end{equation}
  Its connection coefficients belong to $L^2$.
  There are orthonormal parallel frames on the entire interval.
  \item The geodesic velocity has a continuous representative in
  these fibers, satisfies $\mathcal Dv=0$, and has norm equal to
  the constant metric speed.  At $a$ and $b$ the fibers and all
  specified test gradients agree with their tangent-module
  counterparts almost surely.
 \end{enumerate}
\end{proposition}

\begin{proof}
 \textbf{Fibers and connection.}
 Lemmas~\ref{frames:path-avoidance} and
 \ref{frames:regular-paths} ensure that, along almost every entire
 path, the common Gram representatives have rank $n$ at every
 time.  Sobolev differentiation of their entries gives
 \begin{equation}\label{frames:gram-derivative}
  \frac{d}{dt}\langle e_i,e_j\rangle
       =\langle\nabla_v e_i,e_j\rangle
                     +\langle e_i,\nabla_v e_j\rangle.
 \end{equation}
 All terms are in $L^2$ in time along almost every path, since
 the fields are bounded, their covariant derivatives are in
 $L^2(m)$, and speed and compression are bounded.
 Countability permits these assertions simultaneously for every
 pair of specified fields.

 One concrete definition of a fiber is the linear span of the
 symbols $e_j$ modulo the kernel of their Gram form.  All finite
 Gram matrices have rank at most $n$, and some $n$-tuple is
 positive definite, so this fiber has dimension $n$.
 Continuity of Gram entries makes a positive-definite $n$-tuple
 a frame on an open time interval.  Compactness of $[a,b]$
 gives a finite subdivision such that each closed piece lies
 in a frame interval and its inverse Gram matrix is bounded.
 Frame changes have $W^{1,2}$ coefficients, as follows by
 multiplying cross-Gram entries by the inverse Gram matrix.

 At almost every time these fibers are the original tangent
 fibers and the frame spans the tangent space.
 Thus $\nabla_v e_i$ has $L^2$ coefficients in a frame, and
 \eqref{frames:curve-derivative} is meaningful.
 To verify that it is independent of a representation, suppose
 $W=\sum_i a_ie_i$ is zero.  Its pairing with each frame vector
 is identically zero.  Differentiate that pairing using
 \eqref{frames:gram-derivative}; the term pairing $W$ with the
 derivative of the frame vector is zero almost everywhere.
 It follows that the right-hand side of
 \eqref{frames:curve-derivative} pairs to zero against a basis
 almost everywhere and hence vanishes.  The same calculation
 proves metric compatibility.
 In each frame the parallel equation is a linear ODE with
 $L^2$, hence $L^1$, coefficients.  Its absolutely continuous
 solutions exist uniquely.  Metric compatibility preserves
 their pairings, and gluing at the subdivision endpoints
 constructs an orthonormal parallel frame on $[a,b]$.

 \textbf{Parallel velocity.}
 Put $T=b-a$.  Apply \cite[Theorem~5.13(i)]{GigliTamanini2021}
 to the bounded-support optimal geodesic test plan after the affine
 reparameterization $\widetilde\gamma_\tau=\gamma_{a+T\tau}$.
 The functions $g_i$ are in $H^{2,2}$, as required there.
 If $\nabla\widetilde\varphi_\tau$ is its Kantorovich velocity,
 the first differentiation identity identifies the actual velocity
 in the original parameter as
 $v_t=T^{-1}\nabla\widetilde\varphi_{(t-a)/T}$ along the plan
 almost everywhere.  The second differentiation identity has a
 factor $T^2$ on both sides, giving, for $e_i=\nabla g_i$,
 \begin{equation}\label{frames:velocity-pairing}
  \frac{d}{dt}\langle e_i,v\rangle
       =\Hess g_i(v,v)=\langle\nabla_v e_i,v\rangle.
 \end{equation}
 Its identity in $L^2(\pi)$ yields the pathwise distributional
 identity by Fubini; its right-hand side is in
 $L^2(dt\,d\pi)$ by bounded speed and compression.
 Equivalently one may apply the theorem to boundedly weighted
 restrictions of the plan, which remain optimal.
 Choose jointly measurable representatives
 $u_i(t,\gamma)=\langle e_i,v\rangle(\gamma_t)$ and
 $f_i(t,\gamma)=\Hess g_i(v,v)(\gamma_t)$ in $L^2(dt\,d\pi)$.
 Their continuous scalar representatives can be specified by
 \begin{equation}\label{frames:measurable-scalar-trace}
  u_i^*(t,\gamma)=\frac1T\int_a^b u_i(s,\gamma)\dd s
      +\int_a^t f_i(r,\gamma)\dd r
      -\frac1T\int_a^b(b-r)f_i(r,\gamma)\dd r.
 \end{equation}
 This formula is jointly measurable.  Along almost every path it
 belongs to $W^{1,2}$, has derivative $f_i$, and has the same
 average as $u_i$.  It therefore represents the original scalar
 Sobolev class.  Countability gives these assertions on one conull
 path set.

 In a positive-Gram frame put $c=P^{-1}u^*$, where $u^*$ is the
 vector of the corresponding scalar representatives.  The coefficients
 $c$ are $W^{1,2}$.  Write $\nabla_v e_i=\sum_j C_{ji}e_j$.
 Metric compatibility and \eqref{frames:velocity-pairing} give
 \[
  P'=C^{\mathsf T}P+PC,\qquad u^*=Pc,\qquad
  (u^*)'=C^{\mathsf T}Pc.
 \]
 Differentiation yields $c'+Cc=0$, proving $\mathcal Dv=0$
 from the scalar Sobolev identities.  Its almost-everywhere norm
 is the constant metric speed, so continuity gives the same norm
 at every time.

 \textbf{Measurable endpoint traces.}
 The Gram identities defining the endpoint fibers
 agree with the actual tangent-module identities outside an
 $m$-null set.  Endpoint compression excludes this set.
 Countability of the family makes the identification
 simultaneous for all specified fields.  The endpoint vector can
 also be chosen measurably: enumerate the countable $n$-tuples and,
 at $b$, select the first tuple $I=(i_1,\ldots,i_n)$ with positive
 Gram determinant.  This defines a measurable countable partition
 of the good paths.  In a measurable realization of the tangent
 module, set
 \begin{equation}\label{frames:measurable-endpoint-velocity}
  v_b=\sum_{\alpha=1}^n
      \bigl(P_I(b,\gamma)^{-1}u_I^*(b,\gamma)\bigr)_\alpha
                   e_{i_\alpha}(\gamma_b).
 \end{equation}
 This is measurable.  The chosen minor remains positive on a
 neighborhood of $b$ along that path, so the construction agrees
 with the continuous pathwise velocity just obtained.  All its
 other specified gradient pairings equal $u_j^*(b)$ by continuity,
 and its norm equals the path speed.  The same construction applies
 at $a$.  Thus the later endpoint $L^2(\pi)$ pairings use a
 measurable bounded vector defined by scalar traces, independently
 of any choice of values for the original time-equivalence class.
\end{proof}

\begin{lemma}[Countable trials with prescribed global endpoints]
 \label{frames:trial-density}
 Fix $0<\delta<1$, a compact test gradient $Z$, and a countable
 family in Proposition~\ref{frames:countable-frames} whose gradients
 include $Z$.  There is a countable collection $\mathcal Y_{\delta,Z}$
 of fields that are smooth temporal linear combinations of
 these gradients and satisfy
 \[
  Y_\delta=0,\qquad Y_1=Z
 \]
 as identities of fields on $M$, with the following property.
 For every optimal geodesic plan with the bounds of
 Proposition~\ref{frames:path-calculus} on $[\delta,1]$, along
 almost every path every $W^{1,2}$ field $W$ satisfying
 $W(\delta)=0$ and $W(1)=Z(\gamma_1)$ is a limit of members
 of $\mathcal Y_{\delta,Z}$ in the norm
 \[
  \left(\int_\delta^1
       (|Y-W|^2+|\mathcal DY-\mathcal DW|^2)\dd t\right)^{1/2}.
 \]
 The approximating sequence may depend on the path and on $W$.
\end{lemma}

\begin{proof}
 Choose a smooth scalar $\eta$ with $\eta(\delta)=0$ and
 $\eta(1)=1$.  For every interval with rational endpoints
 compactly contained in $(\delta,1)$, choose a countable
 subset of $C_c^\infty$ dense in its $W^{1,2}_0$ space.
 Let $\mathcal A$ be their union and its rational linear span.
 Define $\mathcal Y_{\delta,Z}$ to consist of
 \[
  \eta(t)Z+\sum_{\ell=1}^k a_\ell(t)e_{j_\ell},
       \qquad a_\ell\in\mathcal A,
 \]
 with arbitrary finite lists of indices.  This is countable
 and has the required global endpoint identities.

 Work on a path satisfying Proposition~\ref{frames:path-calculus}.
 Subtract $\eta Z$ from $W$.  In a global parallel frame the
 resulting field has $W^{1,2}_0$ coefficients.  Cutting it off
 near $\delta$ and $1$ converges in $W^{1,2}$; this is the
 usual zero-trace approximation on an interval.
 A smooth temporal partition of unity then expresses the
 remaining compactly supported field in finitely many of the
 local test-gradient frames.  Refine to intervals with rational
 endpoints whose closures remain in these frame intervals.
 The coefficient functions belong to $W^{1,2}_0$ on the
 corresponding intervals.  Approximate them by elements of
 the chosen countable dense subsets.  Such approximation
 converges uniformly as well as in $W^{1,2}$ on an interval.
 Uniform convergence controls multiplication by the $L^2$
 connection coefficients in \eqref{frames:curve-derivative}.
 The fields therefore converge in the stated norm.
\end{proof}

\begin{remark}\label{frames:pathwise-scope}
 Lemma~\ref{frames:trial-density} is a pathwise approximation
 statement with one countable list of global trials.  Below
 it is used by taking their infimum separately at each
 terminal point.  No simultaneous approximation in an
 integrated space of fields along the plan is asserted.
 For related notions of parallel transport along test plans
 and the distinction between the associated Sobolev spaces,
 see \cite{CGP2022}.
\end{remark}

\subsection{The terminal velocity of radial calibrated paths}
\label{frames:sub-radial}

From now on assume $K=(n-1)\kappa$ and the sectional tensor lower
bound, so that Theorem~\ref{heat:index} applies.

Fix $R_\kappa>0$ such that $R_\kappa\sqrt\kappa<\pi/2$ when
$\kappa>0$; for $\kappa\leq0$ fix any positive $R_\kappa$.
For $p\in M$ write $r=d(p,\cdot)$ and $q=r^2/2$.
The bounded Lipschitz function
\begin{equation}\label{frames:radial-data}
 u(x)=\min\{2R_\kappa r(x),\,2R_\kappa^2\}
\end{equation}
is constant off a compact set and satisfies $Q_1u=q$ on
$B(p,R_\kappa)$, with $p$ the unique minimizer.
Indeed, if $s=d(p,y)<R_\kappa$ and $r(x)<R_\kappa$, then
\[
 u(y)+\tfrac12d(x,y)^2-\tfrac12r(x)^2
       \geq s(2R_\kappa-r(x))+\tfrac12s^2>0
       \quad\text{when }s>0.
\]
For $s\geq R_\kappa$ the value of $u$ alone is larger than
$q(x)$.  Thus every calibrated path ending in this ball starts
at $p$, has speed $r(\gamma_1)$, and satisfies
$r(\gamma_t)=t r(\gamma_1)$.

Choose a compact Lipschitz cutoff $\chi$ equal to one on
$\overline{B(p,R_\kappa)}$.  The common terminal-point
prescription in Theorem~\ref{heat:index} yields a finite plan
$\pi_\chi$ with terminal measure $\chi m$.
Its restriction to $[\delta,1]$ has bounded positions, speed,
and compression, and is optimal for every $\delta>0$.
Optimality also follows directly from calibration and the
Hopf--Lax semigroup identity on subintervals.
Restricting this plan by the condition
$\gamma_1\in B(p,R_\kappa)$ gives a plan with terminal measure
$m|_{B(p,R_\kappa)}$ and the same positive-time bounds.
This restriction will be used for representative arguments;
the weights in the index inequality itself remain Lipschitz.

\begin{lemma}[Terminal gradient trace]\label{frames:terminal-trace}
 For the radial calibrated paths just described, the terminal
 representative supplied by Proposition~\ref{frames:path-calculus}
 satisfies
 \begin{equation}\label{frames:terminal-identity}
  v(1)=\nabla q(\gamma_1)
 \end{equation}
 almost surely.  In particular $|\nabla q|=r$ almost everywhere
 on $B(p,R_\kappa)$.
\end{lemma}

\begin{proof}
 Choose a compactly supported Lipschitz function $\widetilde q$
 equal to $q$ on a neighborhood of $\overline{B(p,R_\kappa)}$,
 and compact test functions $h_j\to\widetilde q$ in $W^{1,2}$.
 Include this deterministic countable sequence in the family
 of Proposition~\ref{frames:countable-frames} before applying
 the pathwise statements.  Test-function density here follows
 from heat regularization and good cutoffs; see \cite{Gigli2018}.

 Let $\pi$ denote the radial restriction above, and fix one
 $\delta<1$.  If $(e_t)_\#\pi\leq C m$ on $[\delta,1]$, then
 for any Lipschitz $e\in W^{1,2}$ and $0<h<1-\delta$, Sobolev
 differentiation and Jensen's inequality give
 \begin{equation}\label{frames:quotient-bound}
  \left\|\frac{e(\gamma_1)-e(\gamma_{1-h})}{h}
                                      \right\|_{L^2(\pi)}^2
   \leq\frac{R_\kappa^2}{h}\int_{1-h}^1
                     \int|\nabla e|^2(\gamma_t)\dd\pi\dd t
   \leq CR_\kappa^2\|\nabla e\|_{L^2(m)}^2.
 \end{equation}
 For each fixed $j$, the continuous representative of
 $\langle\nabla h_j,v\rangle$ gives
 \[
  \frac{h_j(\gamma_1)-h_j(\gamma_{1-h})}{h}
       \longrightarrow\langle\nabla h_j(\gamma_1),v(1)\rangle
       \quad\text{in }L^2(\pi).
 \]
 Dominated convergence applies with bound
 $R_\kappa\operatorname{Lip}(h_j)$; no uniform bound in $j$
 is required.  Terminal compression and $|v(1)|\leq R_\kappa$
 also imply
 \[
  \|\langle\nabla(\widetilde q-h_j)(\gamma_1),v(1)\rangle
                                                \|_{L^2(\pi)}
       \leq R_\kappa\sqrt C\,
                    \|\nabla(\widetilde q-h_j)\|_{L^2(m)}.
 \]
 On the radial paths the quotient for $\widetilde q$ is exactly
 $(1-h/2)r(\gamma_1)^2$.
 Apply \eqref{frames:quotient-bound} to
 $e=\widetilde q-h_j$, let $h\downarrow0$ first, and then let
 $j\to\infty$.  This proves
 \[
  \langle\nabla q(\gamma_1),v(1)\rangle=r(\gamma_1)^2.
 \]
 Locality justifies replacing $\widetilde q$ by $q$ at the
 terminal points.  Since $|\nabla q|\leq r$ and
 $|v(1)|=r(\gamma_1)$, equality in Cauchy--Schwarz proves
 \eqref{frames:terminal-identity}, including the zero-speed
 case.  The terminal measure is $m$ on the ball, giving the
 last assertion.
\end{proof}

\subsection{The countable infimum and the Jacobi trial}

Let $\mathfrak s_\kappa$ be the solution of
\[
 \mathfrak s_\kappa''+\kappa\mathfrak s_\kappa=0,
 \qquad \mathfrak s_\kappa(0)=0,
 \qquad \mathfrak s_\kappa'(0)=1.
\]

\begin{theorem}[Distributional Hessian comparison]
 \label{frames:hessian}
 Let \(n\ge2\) be an integer, let \(\kappa\in\mathbb R\), and let
 \((M,d,m)\), with \(m=\mathcal H^n\), be a full-support
 \(\RCD((n-1)\kappa,n)\) space satisfying
 \eqref{eq:main-curvature} for every global \(X,Y\in\TestV(M)\)
 and every nonnegative \(f\in\Test(M)\).
 Fix \(R_\kappa>0\), with \(R_\kappa\sqrt\kappa<\pi/2\) if
 \(\kappa>0\). For every \(p\in M\), put \(r=d(p,\cdot)\)
 and \(q=r^2/2\). Then for every compact test gradient \(Z\),
 $\Hess q(Z,Z)$ on $B(p,R_\kappa)$ is a signed Radon measure
 whose singular part with respect to $m$ is nonpositive.
 It satisfies
 \begin{equation}\label{frames:hessian-bound}
  \Hess q(Z,Z)
    \leq\left(
       |Z_\parallel|^2+
       r\frac{\mathfrak s_\kappa'(r)}{\mathfrak s_\kappa(r)}
                   |Z_\perp|^2\right)m.
 \end{equation}
 Here $|\nabla q|=r$ almost everywhere, and on $\{r>0\}$
 \[
  Z_\parallel=\frac{\langle Z,\nabla q\rangle}{r^2}\nabla q,
             \qquad Z_\perp=Z-Z_\parallel.
 \]
 At $r=0$ the coefficient on the right-hand side of
 \eqref{frames:hessian-bound} is interpreted as $|Z|^2$.
 The measure inequality holds, in particular, against every
 nonnegative compactly supported Lipschitz weight on the ball.
 The radius $R_\kappa$ is independent of $p$.
\end{theorem}

\begin{proof}
 Use the radial data \eqref{frames:radial-data}.
 For the fixed $Z$, adjoin its potential and the approximations
 in Lemma~\ref{frames:terminal-trace} to the countable family.
 Fix $0<\delta<1$.  For every trial \eqref{heat:trial-class}
 satisfying the global endpoint identities $Y_\delta=0$ and $Y_1=Z$,
 Theorem~\ref{heat:index} gives
 \begin{equation}\label{frames:index-terminal}
  \Hess q(Z,Z)(f)\leq\int_{B(p,R_\kappa)} f(x)E_Y(x)\dd m(x),
 \end{equation}
 where $f\geq0$ is compactly supported and Lipschitz in the ball,
 the common calibrated path ending at $x$ is denoted by
 $\gamma^x$, and
 \begin{equation}\label{frames:trial-energy}
  E_Y(x)=\int_\delta^1
    \bigl(|\partial_tY+\nabla_vY|^2
                   -\kappa(|Y|^2|v|^2-\langle Y,v\rangle^2)
    \bigr)(\gamma^x_t)\dd t.
 \end{equation}
 The common terminal prescription is essential here: the
 right-hand side is integration against a function of $x$,
 without averaging over several paths ending at $x$.
 The velocities agree under reweighting of these paths, since
 their pairings with test gradients are the derivatives of
 the same test functions along the paths.
 Bounded compression of $\pi_\chi$, bounded speed, and the
 $L^2$ covariant derivatives of each fixed trial show that
 $E_Y\in L^1_{\mathrm{loc}}(B(p,R_\kappa),m)$.

 We first extract a measure using one fixed trial, for example
 $Y_t=(t-\delta)Z/(1-\delta)$, and denote its energy by $E_0$.
 The functional
 \[
  f\longmapsto\int E_0 f\dd m-\Hess q(Z,Z)(f)
 \]
 is nonnegative on nonnegative compact Lipschitz functions.
 On a fixed compact subset, comparison with a Lipschitz cutoff
 equal to one there bounds its absolute value by a constant
 times $\|f\|_\infty$.  Density in $C_c$ and the Riesz
 representation theorem therefore give a positive Radon
 measure $\nu$ representing it.  Thus
 \[
  \mu:=\Hess q(Z,Z)=E_0m-\nu
 \]
 is a signed Radon measure, with nonpositive singular part.
 Its representation of the fixed functional is unique, so
 the measure does not depend on the initially chosen trial
 or on $\delta$.
 Write $\mu=h_Zm+\mu^{\mathrm s}$ for its Lebesgue decomposition.

 Apply \eqref{frames:index-terminal} to each member of the
 countable collection $\mathcal Y_{\delta,Z}$ from
 Lemma~\ref{frames:trial-density}.  Each inequality says that
 $E_Ym-\mu$ is a positive Radon measure.  Consequently
 \begin{equation}\label{frames:countable-infimum}
  h_Z(x)\leq\inf_{Y\in\mathcal Y_{\delta,Z}}E_Y(x)
                                      \quad\text{for }m\text{-a.e. }x.
 \end{equation}
 The null set can be chosen simultaneously because the list
 is countable.  The same is true for the pathwise calculus
 under $\pi_\chi$ and its radial restriction.

 On such a path, $v$ is parallel and has norm $r=r(x)$.
 Lemma~\ref{frames:terminal-trace} identifies its terminal
 direction with $\nabla q(x)$.  Parallel translate the two
 components $Z_\parallel(x)$ and $Z_\perp(x)$ to obtain
 parallel fields $P_\parallel$ and $P_\perp$ on $[\delta,1]$.
 For $r>0$ take the pathwise field
 \[
  W_t=\frac{t-\delta}{1-\delta}P_\parallel(t)
       +\frac{\mathfrak s_\kappa((t-\delta)r)}
              {\mathfrak s_\kappa((1-\delta)r)}P_\perp(t).
 \]
 Its initial value is zero and its terminal value is $Z(x)$.
 The denominator is positive by the choice of $R_\kappa$.
 At $r=0$ use the linear coefficient on the whole field.
 Lemma~\ref{frames:trial-density} approximates $W$ in the
 pathwise norm by members of $\mathcal Y_{\delta,Z}$.
 The energy \eqref{frames:trial-energy} is continuous in this
 norm because $v$ is bounded.  Hence the infimum in
 \eqref{frames:countable-infimum} is bounded above by the
 energy of $W$.  Only this pathwise upper bound on a countable
 infimum is needed; no measurable selection of the fields
 $W$ or their approximating sequences is required.

 The parallel component contributes
 $|Z_\parallel|^2/(1-\delta)$.  If
 $j(t)=\mathfrak s_\kappa((t-\delta)r)/
             \mathfrak s_\kappa((1-\delta)r)$, then
 $j''+\kappa r^2j=0$, $j(\delta)=0$, and $j(1)=1$.
 Thus the transverse contribution is
 \[
  |Z_\perp|^2\int_\delta^1(j'^2-\kappa r^2j^2)\dd t
   =|Z_\perp|^2[jj']_\delta^1
   =r\frac{\mathfrak s_\kappa'((1-\delta)r)}
               {\mathfrak s_\kappa((1-\delta)r)}|Z_\perp|^2.
 \]
 We have proved
 \[
  h_Z\leq\frac{|Z_\parallel|^2}{1-\delta}
       +r\frac{\mathfrak s_\kappa'((1-\delta)r)}
                {\mathfrak s_\kappa((1-\delta)r)}|Z_\perp|^2
 \]
 almost everywhere, with the continuous zero-speed value
 $|Z|^2/(1-\delta)$.
 Choose a countable sequence $\delta\downarrow0$.
 The same density $h_Z$ occurs for every $\delta$, so outside
 their common null set we may pass to the limit pointwise.
 Together with $\mu^{\mathrm s}\leq0$, this gives
 \eqref{frames:hessian-bound} as a measure inequality and
 therefore against every indicated Lipschitz weight.
\end{proof}
